\documentclass[11pt,a4paper]{article}
\usepackage[utf8]{inputenc}
\usepackage[T1]{fontenc}
\usepackage[spanish]{babel}
\usepackage{hyperref}
\usepackage{geometry}
\geometry{margin=2.5cm}
\usepackage{parskip}

\title{Un alma crece}
\author{Alejandro Valenzuela Rivera\\\small Universidad Autónoma de Nuevo León}
\date{\today}
\begin{document}
\maketitle

\begin{quote}
\textbf{Sesgo declarado.} Este trabajo se sostiene desde el sesgo declarado de su autor, y no finge neutralidad. Lo que afirma se puede tumbar, y ese es el punto: el falsificador está escrito, no escondido.
\end{quote}

\section{La tesis}
\textbf{\textquestiondown{}Puede un alma crecer sin dejar de ser ella?}

\section{Por qué se sostiene}
\textbf{Porque su contexto refutable está escrito, y se puede tumbar sin que esta pregunta caiga.} \textbf{Y porque el borde está escrito: si el criterio cambia, no es la misma alma.}

\section{De dónde viene}
\textbf{De `Tesis/ILUMINAR.md`, y de `base/alma en blanco/MUDAR.vd`.} \textbf{Y de la corrección de la autoridad: «todo muta siempre, se rehace y deshace, y por eso tiene nombre en ese momento, y hasta la tarea cambia».}

\section{Lo que queda abierto}
\textbf{\textquestiondown{}Cómo se distingue un alma que creció de un alma que se sustituyó sin decirlo?}

\section{El caso}
\textbf{Alita, de \emph{Battle Angel Alita} (Yukito Kishiro; Shueisha, 1990): un cuerpo que cambia y un criterio que no.} \textbf{Elige contra su propia misión, y por eso sigue siendo ella.}

\section{Método}
\textbf{\textquestiondown{}La falta está nombrada, y se puede señalar en un archivo? \textquestiondown{}Se buscó la pieza existente antes de crear? \textquestiondown{}Se puede escribir `Criterio: igual`? \textquestiondown{}Se fundió la piel que dejó de servir, con su `Vacío:`?}

\section{Límites, y qué la tumba}
\textbf{Una piel nueva sin falta nombrada; una muda sin término; una muda donde el criterio cambió y no se declaró; una piel fundida sin `Vacío:`; una muda que no corrió nunca; una muda sin autoridad.}

\section{Estado de la tesis}
\textbf{escrita.} \textbf{Falta el caso que la sostenga con una corrida real: hoy el mecanismo está, y el ciclo no ha dado una vuelta completa con autoridad.} ---

\section*{Cómo se leyó esta tesis}
Esto se escribió en el repositorio man-made-empathy, y la tesis está completa en vision del autor/tesis/man made empathy/un alma crece/. La máquina que la compuso no agregó nada que no estuviera escrito: si algo de esto es falso, se dice en el mismo sitio, y se corrige con fecha.
\end{document}
